% @see : https://coopmaths.fr/alea/?uuid=f1208&id=2N30-flash3&n=1&d=10&qcm=0&cd=1&tip=0&alea=IN3l&uuid=d309b&id=2N30-4&n=1&d=10&s=false&s2=10&s3=true&s4=true&cd=0&tip=0&alea=VMzx&uuid=b51bd&id=2N31-1&n=2&d=10&s=1-1-1-0*1111&cd=1&tip=0&alea=omYR&uuid=6575c&id=2N32-1&n=2&d=10&s=4-6&s2=true&cd=0&tip=0&alea=JbG5&uuid=1894f&id=2N31-flash4&alea=kUaf&uuid=29919&id=2N31-2&n=2&d=10&s=1-3-4&s2=true&s3=false&s4=true&s5=1&cd=1&tip=0&alea=Ook5&uuid=88281&id=2N32-flash1&alea=6Kov&uuid=1fc02&id=2N32-2&n=2&d=10&s=4&cd=1&tip=0&alea=T4pN&uuid=67f73&id=2N33-3&alea=ekOk&uuid=7ba8b&id=2N33-2&n=1&d=10&s=8&cd=1&tip=0&alea=6PxW&v=eleve
\documentclass[a4paper,11pt,fleqn]{article}

\usepackage{tasks}
%%% COULEURS %%%
\usepackage[table,svgnames]{xcolor}
\definecolor{nombres}{cmyk}{0,.8,.95,0}
\definecolor{gestion}{cmyk}{.75,1,.11,.12}
\definecolor{gestionbis}{cmyk}{.75,1,.11,.12}
\definecolor{grandeurs}{cmyk}{.02,.44,1,0}
\definecolor{geo}{cmyk}{.62,.1,0,0}
\definecolor{algo}{cmyk}{.69,.02,.36,0}
\definecolor{correction}{cmyk}{.63,.23,.93,.06}
\arrayrulecolor{couleur_theme} % Couleur des filets des tableaux

\usepackage[most]{tcolorbox} % Supprimé à cause de conflits avec ProfCollege
% Découpe des cadres d'exercice. Le style est redéfini par la sortie LaTeX
% juste avant un exercice à garder d'un seul tenant, puis rétabli :
%   \tcbset{mathaleaexo/.style={unbreakable}}
% Passer par un style plutôt que par la clé « breakable » directement évite
% de changer la découpe des autres cadres du document (ProfCollege...).
\tcbset{mathaleaexo/.style={breakable}}
\tcbuselibrary{documentation}

%%%%%%%% Modification paramétrage pour footnotes
% Le paquet semble déjà appelé en amont, curieux qu'il n'y ait pas eu de clash
% avec l'appel ci-dessous mais seulement quand j'ai essayé de l'inclure avec l'option lualatex

% \usepackage{hyperref}
% Parametrages
\hypersetup{
    colorlinks=true,% On active la couleur pour les liens. Couleur par défaut rouge
    linkcolor=black,% On définit la couleur pour les liens internes
    % filecolor=magenta,% On définit la couleur pour les liens vers les fichiers locaux      
    urlcolor=black,% On définit la couleur pour les liens vers des sites web
    % pdftitle={Puissance Quatre},% On définit un titre pour le document pdf
    % pdfpagemode=FullScreen,% On fixe l'affichage par défaut à plein écran
    }
% Pour avoir les numérotations arabic y compris dans les environnements tcolorbox
\renewcommand{\thempfootnote}{\arabic{mpfootnote}}
%%%%%%%%

\usepackage{multicol} 
\usepackage{calc}
\usepackage{enumitem}
\usepackage{graphicx}
\usepackage{tabularx}
\usepackage{tablvar}

\setlength{\parindent}{0mm}
\renewcommand{\arraystretch}{1.5}
\renewcommand{\labelenumi}{\textbf{\theenumi{}.}}
\renewcommand{\labelenumii}{\textbf{\theenumii{}.}}
\setlength{\fboxsep}{3mm}

\setlength{\headheight}{14.5pt}

\setlength{\premulticols}{0pt}

\newcounter{ExoMA}

\newlength{\parindentMA}%
\setlength{\parindentMA}{\parindent}
\addtolength{\parindentMA}{30pt}

\usepackage{simplekv}

\setKVdefault[Theme]{Coopmath,Classiques=false,CANs=false}
\defKV[Theme]{Classique=\setKV[Theme]{Coopmath=false}\setKV[Theme]{Classiques}}
\defKV[Theme]{CAN=\setKV[Theme]{Coopmath=false}\setKV[Theme]{CANs}}

\usepackage{fancyhdr}
\fancyhead[C]{}
\fancyfoot{}

\def\bla{}

\NewDocumentCommand\Theme{o m m m m}{%
  \useKVdefault[Theme]%
  \setKV[Theme]{#1}%
  \ifboolKV[Theme]{CANs}{%
    \usepackage[left=1.5cm,right=1.5cm,top=2cm,bottom=2cm]{geometry}%
    \fancypagestyle{premierePage}{%
      \fancyhead[C]{\textsc{\ifx\bla#2\bla Course aux nombres\else#2\fi}}%
      \fancyfoot[R]{\scriptsize Coopmaths.fr -- CC-BY-SA}%
      \setlength{\headheight}{14.5pt}%
      \NewDocumentCommand\dureeCan{}{\ifx\bla#4\bla 9 minutes\else #4\fi}
      \NewDocumentCommand\titreSujetCan{}{\ifx\bla#3\bla\else\textbf{#3}\fi}
    }%
    \pagestyle{premierePage}
    \colorlet{couleur_theme}{black}%
    \colorlet{couleur_numerotation}{couleur_theme}%
    \let\EXO\EXOCan\let\endEXO\endEXOCan%
  }{%
    \renewcommand{\headrulewidth}{0pt}
    \renewcommand{\footrulewidth}{0pt}
    \pagestyle{fancy}
    \ifboolKV[Theme]{Classiques}{%
      \usepackage[left=0.65cm,right=0.65cm,top=2cm,bottom=1.5cm]{geometry}
      \let\EXO\EXOlibre\let\endEXO\endEXOlibre%
      \colorlet{couleur_theme}{black}%
      \colorlet{couleur_numerotation}{couleur_theme}%
      \fancyhead[C]{\bfseries #3}
      \fancyhead[L]{\bfseries #4}
      \fancyfoot[R]{#5}
    }{%
      \usepackage[left=1.5cm,right=1.5cm,top=4cm,bottom=2cm]{geometry}
      \theme{#2}{#3}{#4}{#5}%
      \let\EXO\EXOcoop\let\endEXO\endEXOcoop%
    }%
  }%
}%

\NewDocumentEnvironment{EXOCan}{m m}{%
}{}

\NewDocumentEnvironment{EXOcoop}{m m}{%
  \begin{tcolorbox}[%
    enhanced,
    mathaleaexo,
    colback=white,
    colframe=white,
    coltitle=black,
    title=\IfNoValueTF{#1}{}{\textmd{#1}},%
    overlay unbroken and first={%
      \IfNoValueTF{#2}{}{%
        \node[
        fill=white,
        anchor=east,
        xshift=10pt,
        text=black
        ]
        at (frame.north east)
        {\scriptsize #2};
      }
      \node[anchor=west,xshift=-20pt,yshift=-15pt] 
      at (frame.north west) {\stepcounter{ExoMA}\numb{\arabic{ExoMA}}};
    }
  ]
}{%
  \end{tcolorbox}
}

\NewDocumentEnvironment{EXOlibre}{m m}{%
  \begin{tcolorbox}[%
    enhanced,
    mathaleaexo,
    colback=white,
    colframe=white,
    coltitle=black,
    title=\stepcounter{ExoMA}\textbf{Exercice \arabic{ExoMA}},%
  ]
  \IfNoValueTF{#1}{}{#1\par}%
}{%
  \end{tcolorbox}%
}


%%% MISE EN PAGE %%%
\usepackage{setspace}
\usepackage{pgf}%
\usepackage{pgfplots}
\pgfplotsset{compat=1.18}
\usetikzlibrary{babel,arrows, arrows.meta,calc,fit,patterns,plotmarks,shapes.geometric,shapes.misc,shapes.symbols,shapes.arrows,shapes.callouts, shapes.multipart, shapes.gates.logic.US,shapes.gates.logic.IEC, er, automata,backgrounds,chains,topaths,trees,petri,mindmap,matrix, calendar, folding,fadings,through,positioning,scopes,decorations.fractals,decorations.shapes,decorations.text,decorations.pathmorphing,decorations.pathreplacing,decorations.footprints,decorations.markings,shadows}

\spaceskip=2\fontdimen2\font plus 3\fontdimen3\font minus3\fontdimen4\font\relax %Pour doubler l'espace entre les mots
\newcommand\numb[1]{ % Dessin autour du numéro d'exercice
  \begin{tikzpicture}[overlay,scale=.8]%,yshift=-.3cm
    \draw[fill=couleur_numerotation,couleur_numerotation](-.4,-0.1)rectangle($(-0.4,.8)+(2em,0)$);
    \draw (-.4,.8) node[white,anchor=north west,inner sep=0.5pt]{\bfseries EX}; 
    \draw[line width=.05cm,couleur_numerotation,fill=white] (0,0)--(.5,.5)--(1,0)--(.5,-.5)--cycle;
    \draw (0.5,0) node[anchor=center,couleur_numerotation]{\large\bfseries#1};
  \end{tikzpicture}
}

% echelle pour le dé
\def\globalscale {0.04}
% abscisse initiale pour les chevrons
\def\xini {3}

\newcommand\theme[4]{%
  \fancyhead[C]{%
    % Tracé du dé
    \begin{tikzpicture}[y=0.80pt, x=0.80pt, yscale=-\globalscale, xscale=\globalscale,remember picture, overlay,transform canvas={xshift=-0.5\linewidth,yshift=2.7cm},fill=couleur_theme]%,xshift=17cm,yshift=9.5cm
      %%%% Arc supérieur gauche%%%%
      \path[fill](523,1424)..controls(474,1413)and(404,1372)..(362,1333)..controls(322,1295)and(313,1272)..(331,1254)..controls(348,1236)and(369,1245)..(410,1283)..controls(458,1328)and(517,1356)..(575,1362)..controls(635,1368)and(646,1375)..(643,1404)..controls(641,1428)and(641,1428)..(596,1430)..controls(571,1431)and(538,1428)..(523,1424)--cycle;
      %%%% Dé face supérieur%%%%
      \path[fill](512,1272)..controls(490,1260)and(195,878)..(195,861)..controls(195,854)and(198,846)..(202,843)..controls(210,838)and(677,772)..(707,772)..controls(720,772)and(737,781)..(753,796)..controls(792,833)and(1057,1179)..(1057,1193)..controls(1057,1200)and(1053,1209)..(1048,1212)..controls(1038,1220)and(590,1283)..(551,1282)..controls(539,1282)and(521,1278)..(512,1272)--cycle;
      %%%% Dé faces gauche et droite%%%%
      \path[fill](1061,1167)..controls(1050,1158)and(978,1068)..(900,967)..controls(792,829)and(756,777)..(753,756)--(748,729)--(724,745)..controls(704,759)and(660,767)..(456,794)..controls(322,813)and(207,825)..(200,822)..controls(193,820)and(187,812)..(187,804)..controls(188,797)and(229,688)..(279,563)..controls(349,390)and(376,331)..(391,320)..controls(406,309)and(462,299)..(649,273)..controls(780,254)and(897,240)..(907,241)..controls(918,243)and(927,249)..(928,256)..controls(930,264)and(912,315)..(889,372)..controls(866,429)and(848,476)..(849,477)..controls(851,479)and(872,432)..(897,373)..controls(936,276)and(942,266)..(960,266)..controls(975,266)and(999,292)..(1089,408)..controls(1281,654)and(1290,666)..(1290,691)..controls(1290,720)and(1104,1175)..(1090,1180)..controls(1085,1182)and (1071,1176)..(1061,1167)--cycle;
      %%%% Arc inférieur bas%%%%
      \path[fill](1329,861)..controls(1316,848)and(1317,844)..(1339,788)..controls(1364,726)and(1367,654)..(1347,591)..controls(1330,539)and(1338,522)..(1375,526)..controls(1395,528)and(1400,533)..(1412,566)..controls(1432,624)and(1426,760)..(1401,821)..controls(1386,861)and(1380,866)..(1361,868)..controls(1348,870)and(1334,866)..(1329,861)--cycle;
      %%%% Arc inférieur gauche%%%%
      \path[fill](196,373)..controls(181,358)and(186,335)..(213,294)..controls(252,237)and(304,190)..(363,161)..controls(435,124)and(472,127)..(472,170)..controls(472,183)and(462,192)..(414,213)..controls(350,243)and(303,283)..(264,343)..controls(239,383)and(216,393)..(196,373)--cycle;
    \end{tikzpicture}
    \begin{tikzpicture}[line cap=round,line join=round,remember picture, overlay, shift={(current page.north west)},yshift=-8.5cm]
      \fill[fill=couleur_theme] (0,5) rectangle (21,6);
      \fill[fill=couleur_theme] (\xini,6)--(\xini+1.5,6)--(\xini+2.5,7)--(\xini+1.5,8)--(\xini,8)--(\xini+1,7)-- cycle;
      \fill[fill=couleur_theme] (\xini+2,6)--(\xini+2.5,6)--(\xini+3.5,7)--(\xini+2.5,8)--(\xini+2,8)--(\xini+3,7)-- cycle;  
      \fill[fill=couleur_theme] (\xini+3,6)--(\xini+3.5,6)--(\xini+4.5,7)--(\xini+3.5,8)--(\xini+3,8)--(\xini+4,7)-- cycle;   
      \node[color=white] at (10.5,5.5) {\LARGE \bfseries{ \MakeUppercase{#4}}};
    \end{tikzpicture}
    \begin{tikzpicture}[remember picture,overlay]
      \node[anchor=north east,inner sep=0pt] at ($(current page.north east)+(0,-.8cm)$) {};
      \node[anchor=west, fill=white, inner sep = 0pt] at ($(current page.north west)+(0.2,-2.3cm)$) {\footnotesize \bfseries{MathALÉA}};
    \end{tikzpicture}
    \begin{tikzpicture}[remember picture,overlay]
      \node[anchor=north east,inner sep=0pt] at ($(current page.north east)+(0,-.8cm)$) {};
      \node[anchor=east, fill=white] at ($(current page.north east)+(-2,-1.5cm)$) {\Huge \textcolor{couleur_theme}{\bfseries{\#}} \bfseries{#2} \textcolor{couleur_theme}{\bfseries \MakeUppercase{#3}}};
    \end{tikzpicture}
  }%
  \fancyfoot[R]{%
    \begin{tikzpicture}[remember picture,overlay]
      \node[anchor=south east] at ($(current page.south east)+(-2,0.25cm)$) {\scriptsize {\bfseries \href{https://coopmaths.fr/}{Coopmaths.fr} -- \href{http://creativecommons.fr/licences/}{CC-BY-SA}}};
    \end{tikzpicture}
    \begin{tikzpicture}[line cap=round,line join=round,remember picture, overlay,xscale=0.5,yscale=0.5, shift={(current page.south west)},xshift=35.7cm,yshift=-6cm]
      \fill[fill=couleur_theme] (\xini,6)--(\xini+1.5,6)--(\xini+2.5,7)--(\xini+1.5,8)--(\xini,8)--(\xini+1,7)-- cycle;
      \fill[fill=couleur_theme] (\xini+2,6)--(\xini+2.5,6)--(\xini+3.5,7)--(\xini+2.5,8)--(\xini+2,8)--(\xini+3,7)-- cycle;  
      \fill[fill=couleur_theme] (\xini+3,6)--(\xini+3.5,6)--(\xini+4.5,7)--(\xini+3.5,8)--(\xini+3,8)--(\xini+4,7)-- cycle;  
    \end{tikzpicture}
  }
  \fancyfoot[C]{}
  \colorlet{couleur_theme}{#1}
  \colorlet{couleur_numerotation}{couleur_theme}
  \def\iconeobjectif{icone-objectif-#1}
  \def\urliconeomethode{icone-methode-#1}
}%

\newcommand*\tikzfootMA{%
  \ifboolKV[Theme]{CANs}{
    \begin{tikzpicture}[remember picture,overlay,shift=(current page.north west)]
      \begin{scope}[x=\textwidth-\oddsidemargin,y=\textheight+5pt]
        \draw[correction,line width=4pt,dashed,dash pattern= on 10pt off 10pt,shift={(-\oddsidemargin,\topmargin)}] (0,0) rectangle (1,-1);
      \end{scope}
    \end{tikzpicture} 
  }{%
  \ifboolKV[Theme]{Classiques}{%
  \begin{tikzpicture}[remember picture,overlay,shift=(current page.south west)]
    \begin{scope}[x={(current page.south east)},y={(current page.north west)}]
      \draw[correction,line width=4pt,dashed,dash pattern= on 10pt off 10pt] ($(0,0)+(5mm,1cm)$) rectangle ($(1,1)+(-5mm,-1.75cm)$);
    \end{scope}
  \end{tikzpicture} 
  }{%
  \begin{tikzpicture}[remember picture,overlay,shift=(current page.south west)]
    \begin{scope}[x={(current page.south east)},y={(current page.north west)}]
      \draw[correction,line width=4pt,dashed,dash pattern= on 10pt off 10pt] ($(0,0)+(5mm,1.5cm)$) rectangle ($(1,1)+(-5mm,-3.75cm)$);
    \end{scope}
  \end{tikzpicture} 
  }%
  }
}

\newenvironment{Correction}{%
  \setcounter{ExoMA}{0}%
  \cfoot{\tikzfootMA}
}{}%

\newcommand\version[1]{
  \fancyhead[R]{
    \begin{tikzpicture}[remember picture,overlay]
    \node[anchor=north east,inner sep=0pt] at ($(current page.north east)+(-.5,-.5cm)$) {\large \textcolor{couleur_theme}{\bfseries V#1}};
    \end{tikzpicture}
  }
}

\usepackage[french]{babel}
\usepackage{icomma}
% loadPackagesFromContent

\newcounter{customfootnote}
\newcounter{tempreset}%
\newcounter{temp}

\makeatletter
\newcommand{\anote}[1]{%
  \refstepcounter{customfootnote}%
  \textsuperscript{\thecustomfootnote}%
  \expandafter\gdef\csname custom@footnote@\thecustomfootnote\endcsname{#1}%
}

\newcommand{\printcustomnotes}{%
  \par\noindent\rule{\linewidth}{0.4pt}\par
  \setcounter{temp}{1}%
  \whileboolexpr{
    test {\ifcsdef{custom@footnote@\thetemp}}
  }{%
    \noindent\textsuperscript{\thetemp}~\csname custom@footnote@\thetemp\endcsname\par
    \stepcounter{temp}%
  }%
}

\newcommand{\resetcustomnotes}{%
  \setcounter{customfootnote}{0}%
  \setcounter{tempreset}{1}%
  \whileboolexpr{
    test {\ifcsdef{custom@footnote@\thetempreset}}
  }{%
    \global\expandafter\let\csname custom@footnote@\thetempreset\endcsname\@undefined
    \stepcounter{tempreset}%
  }%
}
\makeatother
\usepackage[gen]{eurosym}
\usepackage{cancel}
\usepackage{tikz}
\usetikzlibrary{decorations.pathmorphing}
\usetikzlibrary {decorations.pathreplacing}
\newcommand{\e}{\mathrm{e}}
\usepackage{amsmath}
\usepackage{etoolbox}
\newbool{dys}
\setbool{dys}{false}
% ===== POLICES =====
% Le choix est le même en police standard et en police adaptée aux dys ;
% seuls la taille et l'interligne changent entre les deux.
\newcommand{\choiceFonts}[1]{
\ifstrequal{#1}{helvet}{%
  % Helvetica (police standard historique des fiches)
  \usepackage[scaled=1]{helvet}
}{}
\ifstrequal{#1}{Fira}{%
  % Fira Sans + Fira Math
  % Description : Fira Sans est une police moderne, et Fira Math est une version mathématique compatible qui maintient le style sans-serif.
  % Utilisation : Fonctionne bien avec XeLaTeX et LuaLaTeX.
  \usepackage{fontspec}
  \setmainfont{Fira Sans}
  \setsansfont{Fira Sans}
  \usepackage{unicode-math}
  \setmathfont{Fira Math}
}{}
\ifstrequal{#1}{lmodern}{%
  % Latin Modern Sans
  % Description : Une version sans-serif de la célèbre police Latin Modern, qui est compatible avec mathastext.
  % Utilisation : Fonctionne avec pdfLaTeX, XeLaTeX et LuaLaTeX.
  \usepackage{lmodern}
  \renewcommand{\familydefault}{\sfdefault}
  \usepackage{mathastext}
}{}
\ifstrequal{#1}{tgheros}{%
  % TeX Gyre Heros + mathastext
  % Description : TeX Gyre Heros est une alternative à Helvetica, et le package mathastext permet d'utiliser la même police pour les mathématiques.
  % Utilisation : Fonctionne avec pdfLaTeX, XeLaTeX, ou LuaLaTeX.
  \usepackage{tgheros}
  \renewcommand{\familydefault}{\sfdefault}
  \usepackage{mathastext}
}{}
}
\ifbool{dys}{
% POLICE DYS
\choiceFonts{Fira}

%\usepackage{unicode-math}
%\usepackage{fontspec}
%\setmainfont{TeX Gyre Schola}
%\setsansfont{TeX Gyre Schola}
%\setmathfont{TeX Gyre Schola Math}
%\setmainfont{OpenDyslexic}[Scale=1.0]
%\setsansfont{OpenDyslexic}[Scale=1.0]
%\setmathfont{OpenDyslexic}[Scale=1.0,range=up/{Latin,latin,num}]
\usepackage[fontsize=14]{scrextend}
\usepackage{setspace}
\setstretch{1.7}
% Une valeur d'environ 1.2 à 1.7 est souvent recommandée. Cela permet d'augmenter l'espace entre les lignes tout en offrant un léger espacement entre les mots.
\setlength{\spaceskip}{1.2em}
% Une valeur d'environ 1.2em à 1.5em est couramment conseillée. Cela crée un espace plus ample entre les mots, ce qui peut aider à réduire la fatigue visuelle et à améliorer la fluidité de la lecture.
}{
% POLICE STANDARD
%%% EE (24/04/2026) : Cette modif ci-dessous est nécessaire pour accepter ’ comme apostrophe.
% \usepackage[T1]{fontenc}     % Réservé à pdfLaTeX, à remplacer par fontspec en LuaLaTeX
\choiceFonts{helvet}
\usepackage[fontsize=12]{scrextend}
}

\Theme[Classique]{nombres}{Fracions 2}{Secondes}{Accompagnement Personnalisé}
\begin{document}

\tcbset{mathaleaexo/.style={breakable}}
% @see : https://coopmaths.fr/alea?uuid=f1208&id=2N30-flash3&n=1&d=10&alea=IN3l&cd=1&cols=1
% \begin{EXO}{}{}
% 
% 
% 
%  Rendre  la fraction $\dfrac{9}{24}$ irréductible.
% \tcbset{mathaleaexo/.style={breakable}}
% % @see : https://coopmaths.fr/alea?uuid=d309b&id=2N30-4&n=1&d=10&s=false&s2=10&s3=true&s4=true&alea=VMzx&cd=0&cols=2
% \end{EXO}

\begin{EXO}{}{}
Compléter avec deux nombres entiers \textbf{consécutifs}.

\begin{multicols}{2}
 \begin{minipage}[t]{\linewidth} $\ldots\ldots$$ < \dfrac{33}{6} < $$\ldots\ldots$ \end{minipage}
\end{multicols}
\tcbset{mathaleaexo/.style={breakable}}
% @see : https://coopmaths.fr/alea?uuid=b51bd&id=2N31-1&n=2&d=10&s=1-1-1-0*1111&alea=omYR&cd=1&cols=1
\end{EXO}

\begin{EXO}{}{}
Calculer et donner le résultat sous la forme d’une fraction irréductible.


\begin{tasks}(2)
	\task $1 - \dfrac{1}{6}$
	\task $\dfrac{3}{4} + \dfrac{1}{2}$
\end{tasks}
\tcbset{mathaleaexo/.style={breakable}}
% @see : https://coopmaths.fr/alea?uuid=6575c&id=2N32-1&n=2&d=10&s=4-6&s2=true&alea=JbG5&cd=0&cols=4
\end{EXO}

\begin{EXO}{}{}
Effectuer les calculs suivants en donnant le résultat sous forme d'une fraction.

 \begin{tasks}(3)
	\task $2 \div  \dfrac{4}{7} $
	\task $\dfrac{7}{4} - \dfrac{6}{4} \times \dfrac{4}{9}$ 
    \task $\dfrac{14}{18}\times \dfrac{30}{14}$
 \end{tasks}
\tcbset{mathaleaexo/.style={breakable}}
% @see : https://coopmaths.fr/alea?uuid=1894f&id=2N31-flash4&n=1&d=10&alea=kUaf&cd=1&cols=1
\end{EXO}

\begin{EXO}{}{}
Calculer et donner le résultat sous forme irréductible.


 \begin{tasks}(3)
	\task $44\times\dfrac{5}{66}$
	\task $\dfrac{-21}{-40}\times\dfrac{5}{9}$
    \task $\dfrac{4}{3-\dfrac{1}{5}}$
 \end{tasks}
\tcbset{mathaleaexo/.style={breakable}}
% @see : https://coopmaths.fr/alea?uuid=88281&id=2N32-flash1&n=1&d=10&alea=6Kov&cd=1&cols=1
\end{EXO}

\begin{EXO}{}{}
Calculer et donner le résultat sous la forme d’une fraction irréductible.

\begin{multicols}{2}
\begin{enumerate}
	\item \begin{minipage}[t]{\linewidth} $A=\dfrac{\dfrac{3}{5}+3}{\dfrac{2}{4}-4}$ \end{minipage}
	\item \begin{minipage}[t]{\linewidth} $B=\dfrac{1}{7}\times6+\dfrac{6}{8}\times6$ \end{minipage}
\end{enumerate}
\end{multicols}
\tcbset{mathaleaexo/.style={breakable}}
% @see : https://coopmaths.fr/alea?uuid=67f73&id=2N33-3&n=1&d=10&s=3&alea=ekOk&cd=1&cols=1
\end{EXO}

\begin{EXO}{}{}



 Kevin a couru $\dfrac{3}{7}$ de la distance d'une course à pied la première heure. Ensuite, il a couru $\dfrac{1}{7}$ de la distance totale la deuxième heure.

$\textbf{a)}$ Quelle fraction de la distance totale a-t-il courue ? 

            $\textbf{b)}$ Quelle fraction de la distance totale lui reste-t-il à courir ? 
\tcbset{mathaleaexo/.style={breakable}}
% @see : https://coopmaths.fr/alea?uuid=7ba8b&id=2N33-2&n=1&d=10&s=8&alea=6PxW&cd=1&cols=1
\end{EXO}

\begin{EXO}{}{}
Justifier vos réponses aux problèmes suivants.


 Tania dépense $\dfrac{2}{5}$ de son argent de poche lors du premier jour des soldes.

Le lendemain, elle dépense $\dfrac{5}{12}$ de ce qui lui reste.  

 \textbf {a)}  Quelle fraction de son argent de poche lui reste-il? 

 \textbf {b)}  Sachant qu'elle avait 80\,\euro{}, combien lui reste-t-il? 
\end{EXO}


% \clearpage
% 
% \begin{Correction}
% \begin{EXO}{}{}
% 
%  $\dfrac{9}{24} =\dfrac{3{\color[HTML]{2563a5}\boldsymbol{\times3}} }{8{\color[HTML]{2563a5}\boldsymbol{\times3}}}={\color[HTML]{F15929}\boldsymbol{\dfrac{3}{8}}}$
% \end{EXO}
% 
% \begin{EXO}{}{}
% 
%   $\quad \dfrac{33}{6}=5+\dfrac{3}{6}\quad$ et $\quad5<5+\dfrac{3}{6}<6$ 
% 
% donc $\quad {\color[HTML]{F15929}\boldsymbol{5}} < \dfrac{33}{6} < {\color[HTML]{F15929}\boldsymbol{6}}$.
% \end{EXO}
% 
% \begin{EXO}{}{}
% 
% \begin{enumerate}[label={}]
% 	\item $\begin{aligned}A&=1 - \dfrac{1}{6}\\
% &=\dfrac{6}{6} - \dfrac{1}{6}\\
% &=\dfrac{6-1}{6}\\
% &={\color[HTML]{F15929}\boldsymbol{\dfrac{5}{6}}}
%         \end{aligned}$
% 	\item $\begin{aligned}B&=\dfrac{3}{4} + \dfrac{1}{2}\\
% &=\dfrac{3}{4} + \dfrac{2}{4}\\
% &=\dfrac{3+2}{4}\\
% &={\color[HTML]{F15929}\boldsymbol{\dfrac{5}{4}}}
%         \end{aligned}$
% \end{enumerate}
% \end{EXO}
% 
% \begin{EXO}{}{}
% \begin{multicols}{4}
% \begin{enumerate}[itemsep=1em,label={}]
% 	\item \begin{minipage}[t]{\linewidth} $A = 2 \div  \dfrac{4}{7} $
% 
% $A = 2 \times \dfrac{7}{4}$
% 
% $A = \dfrac{14}{4}$
% 
% $A  =\dfrac{7{\color[HTML]{216D9A}\boldsymbol{\times2}}}{2{\color[HTML]{216D9A}\boldsymbol{\times2}}}={\color[HTML]{F15929}\boldsymbol{\dfrac{7}{2}}}$
% 
%  \end{minipage}
% 	\item \begin{minipage}[t]{\linewidth} $B = \dfrac{7}{4} - \dfrac{6}{4} \times \dfrac{4}{9}$
% 
% $B = \dfrac{7}{4} - \dfrac{24}{36}$
% 
% $B = \dfrac{63}{36} - \dfrac{24}{36}$
% 
% $B = \dfrac{39}{36}$
% 
% $B  =\dfrac{13{\color[HTML]{216D9A}\boldsymbol{\times3}}}{12{\color[HTML]{216D9A}\boldsymbol{\times3}}}={\color[HTML]{F15929}\boldsymbol{\dfrac{13}{12}}}$
% 
%  \end{minipage}
% \end{enumerate}
% \end{multicols}
% \end{EXO}
% 
% \begin{EXO}{}{}
% 
%  Il y a une simplification évidente par $14$, puis on simplifie par $6$ :
% 
%     $\begin{aligned}
%     \dfrac{14}{18}\times \dfrac{30}{14}&=\dfrac{30}{18}\\
%     &=\dfrac{6 \times 5}{6 \times 3}\\
%     &={\color[HTML]{F15929}\boldsymbol{\dfrac{5}{3}}}
%     \end{aligned}$
% \end{EXO}
% 
% \begin{EXO}{}{}
% 
% \begin{enumerate}[itemsep=1em,label={}]
% 	\item $\begin{aligned}A &= 44\times \dfrac{5}{66}\\&=\dfrac{44\times 5}{66}\\&=\dfrac{\cancel{2}\times2\times5\times\cancel{11}}{\cancel{2}\times3\times\cancel{11}}\\&={\color[HTML]{F15929}\boldsymbol{\dfrac{10}{3}}}\end{aligned}$
% 	\item $\begin{aligned}B &= \dfrac{-21}{-40}\times \dfrac{5}{9}\\&=\dfrac{-21\times 5}{-40 \times 9}\\&=\dfrac{(-3)\times7\times 5}{(-2)\times2\times2\times5\times 3\times3}\\&=\dfrac{\cancel{3}\times\cancel{5}\times7}{2\times2\times2\times\cancel{3}\times3\times\cancel{5}}\\&={\color[HTML]{F15929}\boldsymbol{\dfrac{7}{24}}}\end{aligned}$
% \end{enumerate}
% \end{EXO}
% 
% \begin{EXO}{}{}
% 
%  
% $\begin{aligned}
% A &= \dfrac{4}{3-\dfrac{1}{5}} \\
% &= \dfrac{4}{\dfrac{15}{5}-\dfrac{1}{5}} \\
% &= \dfrac{4}{\dfrac{14}{5}} \\
% &= 4 \times \dfrac{5}{14} \\
% &= {\color[HTML]{F15929}\boldsymbol{\dfrac{10}{7}}}
% \end{aligned}$
% \end{EXO}
% 
% \begin{EXO}{}{}
% 
% \begin{enumerate}[itemsep=1em]
% 	\item $\begin{aligned}A&=\dfrac{\dfrac{3}{5}+3}{\dfrac{2}{4}-4}\\&=\dfrac{\dfrac{3}{5}+\dfrac{15}{5}}{\dfrac{2}{4}-\dfrac{16}{4}}&\text{On met au même dénominateur dans le numérateur et le dénominateur.}\\&=\dfrac{\dfrac{3+15}{5}}{\dfrac{2-16}{4}}&\text{On calcule les sommes ou les différences.}\\&=\dfrac{\dfrac{18}{5}}{-\dfrac{14}{4}}\\&=\dfrac{18}{5}\times-\dfrac{2}{7}&\text{Diviser par un réel non nul, c’est multiplier par son inverse.}\\&=-\dfrac{72}{70}\\&={\color[HTML]{F15929}\boldsymbol{-\dfrac{36}{35}}}\end{aligned}$
% 	\item $\begin{aligned}B&={\color{red}\boldsymbol{\dfrac{1}{7}\times6}}+{\color{red}\boldsymbol{\dfrac{6}{8}\times6}}&\text{On repère la priorité des multiplications sur l’addition.}\\&=\dfrac{6}{7}+\dfrac{9}{2}&\text{On effectue les multiplications en priorité.}\\&=\dfrac{12}{14}+\dfrac{63}{14}&\text{On met au même dénominateur.}\\&=\dfrac{12+63}{14}&\text{On calcule la somme ou la différence des numérateurs.}\\&={\color[HTML]{F15929}\boldsymbol{\dfrac{75}{14}}}\end{aligned}$
% \end{enumerate}
% \end{EXO}
% 
% \begin{EXO}{}{}
% 
%  
% \begin{enumerate}[label=\alph*.]
% 	\item Kevin a couru $\dfrac{4}{7}$ de la distance totale.
% 
%  En effet :
% 
% $\begin{aligned}\dfrac{3}{7} + \dfrac{1}{7} &=\dfrac{3+1}{7}\\
%                   &={\color[HTML]{F15929}\boldsymbol{\dfrac{4}{7}}}\end{aligned}$
% 
% \begin{tikzpicture}[scale=0.8]
% \draw[fill=lightgray] (0.0,2.2) rectangle (11.2,3.0) node[pos=0.5] {\textcolor{black}{distance totale}};
% \draw[fill=lightgray] (0.0,1.4) rectangle (4.8,2.2) node[pos=0.5] {\textcolor{black}{$1^{re}$ part}};
% \draw[fill=lightgray] (4.8,1.4) rectangle (6.4,2.2) node[pos=0.5] {\textcolor{black}{$2^e$ part}};
% \draw[fill=lightgray] (0.0,0.6) rectangle (1.6,1.4) node[pos=0.5] {\textcolor{black}{$\phantom{P}$}};
% \draw[fill=lightgray] (1.6,0.6) rectangle (3.2,1.4) node[pos=0.5] {\textcolor{black}{$\phantom{P}$}};
% \draw[fill=lightgray] (3.2,0.6) rectangle (4.8,1.4) node[pos=0.5] {\textcolor{black}{$\phantom{P}$}};
% \draw[fill=lightgray] (4.8,0.6) rectangle (6.4,1.4) node[pos=0.5] {\textcolor{black}{$\phantom{P}$}};
% \draw[decorate,decoration={brace,amplitude=10pt},xshift=0pt,yshift=0pt, draw=black] (4.8,0.6000000000000001) -- node[below=10pt, pos=0.5] {\textcolor{black}{$\dfrac{3}{7}$}} (0.0,0.6000000000000001);
% \draw[decorate,decoration={brace,amplitude=10pt},xshift=0pt,yshift=0pt, draw=black] (6.4,0.6000000000000001) -- node[below=10pt, pos=0.5] {\textcolor{black}{$\dfrac{1}{7}$}} (4.8,0.6000000000000001);
% \end{tikzpicture}
% 
% 	\item Il lui reste $\dfrac{3}{7}$ de la distance totale à courir.
% 
%  En effet :
% 
% $\begin{aligned}1 - \dfrac{4}{7}&=\dfrac{7}{7} - \dfrac{4}{7}\\
%               & = \dfrac{7 - 4}{7}\\
%               &={\color[HTML]{F15929}\boldsymbol{\dfrac{3}{7}}}\end{aligned}$
% 
% \begin{tikzpicture}[scale=0.8]
% \draw[fill=lightgray] (0.0,2.2) rectangle (11.2,3.0) node[pos=0.5] {\textcolor{black}{distance totale}};
% \draw[fill=lightgray] (0.0,0.2) rectangle (6.4,2.2) node[pos=0.5] {\textcolor{black}{$\dfrac{4}{7}$}};
% \draw[fill=lightgray] (6.4,0.2) rectangle (11.2,2.2) node[pos=0.5] {\textcolor{black}{?}};
% \end{tikzpicture}
% \end{enumerate}
% 
% \end{EXO}
% 
% \begin{EXO}{}{}
% 
%  Il s'agit d'un problème multiplicatif. Prendre une fraction d'une fraction revient à multiplier les deux fractions entre elles.
% 
% $1 - \dfrac{2}{5} = \dfrac{3}{5}$ de son argent qui reste après le premier jour.
% 
% $\dfrac{5}{12}\times \dfrac{3}{5} = \dfrac{1}{4} $ de son argent de poche dépensé lors du deuxième jour.
% 
% Finalement, $\dfrac{2}{5}$ de son argent de poche dépensé lors du premier jour et $\dfrac{1}{4}$ de son argent de poche dépensé lors du deuxième jour.
% 
% Réduisons les fractions au même dénominateur :  $\dfrac{2}{5}$ $= \dfrac{8}{20}$ et $\dfrac{1}{4}$ $= \dfrac{5}{20}$.
% 
% Calculons la fraction d'argent qui  lui reste :
% 
% $1-\dfrac{2}{5}-\dfrac{1}{4} = \dfrac{20}{20}-\dfrac{8}{20}-\dfrac{5}{20} = \dfrac{20-8-5}{20} = \dfrac{7}{20}$
% 
% Conclusion : ${\color[HTML]{F15929}\boldsymbol{\dfrac{7}{20}}}$ de son argent de poche qui lui reste après les soldes.
% 
% Sachant qu'il y avait  $80$\,\euro{}, on peut calculer :
% 
% $\dfrac{7}{20}\times 80 = 28 $  \,\euro{}~ lui reste.
% 
% Conclusion : ${\color[HTML]{F15929}\boldsymbol{28}}$\,\euro{}~lui reste.
% \end{EXO}
% 
% \clearpage
% \end{Correction}
\end{document}

% Local Variables:
% TeX-engine: luatex
% End: